\chapter{Barron norm and infinite-width networks}
\label{chap:barron}

\section{Introduction}

This chapter formalizes Chapter~3 of the deep learning theory notes (Telgarsky~2021),
which develops two influential ideas.

\begin{enumerate}
  \item \textbf{Infinite-width networks.}  Instead of a finite parameter tuple, a shallow
    network is described by a \emph{signed measure} $\nu$ over weight space.  The network
    evaluates as $x \mapsto \int \sigma(w^\top x)\,d\nu(w)$, and its \emph{mass} is the
    total variation $|\nu|(\mathbb{R}^p)$.  Infinite-width networks simplify analysis,
    and certain important phenomena — such as staying near random initialization — are
    better captured at this level of abstraction before sampling down to finite width.

  \item \textbf{Barron norm and representation theorem.}  Barron~(1993) identified a class
    of functions — those with a finite ``Barron norm'' defined via the Fourier gradient —
    that can be represented \emph{exactly} as infinite-width threshold networks.  Combining
    this with a sampling argument (Maurey's lemma; Pisier~1980) gives a quantitative bound
    on how many nodes are needed to approximate such a function to within $\varepsilon$
    in $L^2(P)$: roughly $C^2/\varepsilon^2$, where $C$ is the Barron norm.
    Crucially, this complexity is dimension-free, circumventing the curse of dimension for
    functions in the Barron class.
\end{enumerate}

\paragraph{Relation to the broader LML project.}
The Lean formalizations developed in this chapter extend the
\texttt{Approximation} module of \texttt{LeanMachineLearning.Optimization}.
The four new files are:

\begin{itemize}
  \item \texttt{InfiniteWidth.lean} — Definition~\ref{def:infiniteWidthNetwork} and
    Proposition~\ref{prop:univariateIntegralRep}.
  \item \texttt{BarronNorm.lean} — Definition~\ref{def:barronNorm} and
    Theorem~\ref{thm:barronRepresentation}.
  \item \texttt{Sampling.lean} — Lemma~\ref{lem:maureySampling},
    Lemma~\ref{lem:maureySignedMeasure}, and
    Corollary~\ref{cor:barronSamplingBound}.
\end{itemize}

\paragraph{Notation.}
Throughout this chapter, $d \in \mathbb{N}$ is the input dimension,
$\sigma_\theta$ denotes the threshold activation $z \mapsto \mathbf{1}[z \ge 0]$,
$\widehat{f}$ denotes the Fourier transform of $f$ with convention
$\widehat{f}(w) = \int \exp(-2\pi i\, w^\top x)\, f(x)\, dx$,
and $P$ denotes a probability measure on inputs $x \in \mathbb{R}^d$.
We write $\|f\|_{L_2(P)}^2 = \int f(x)^2\, dP(x)$.

% -------------------------------------------------------------------------
\section{Infinite-width shallow networks}
\label{sec:infinite-width}

\subsection{The infinite-width model}

\begin{definition}[Infinite-width shallow network]\label{def:infiniteWidthNetwork}
  \lean{Approximation.InfiniteWidth.InfiniteWidthNetwork}
  \uses{def:thresholdActivation, def:reluActivation}
An \emph{infinite-width shallow network} with activation $\sigma$ and weight space
$\mathbb{R}^p$ is a pair $(\sigma, \nu)$ where $\nu$ is a signed Borel measure on
$\mathbb{R}^p$.
The network evaluates as
\[
  x \;\mapsto\; \int_{\mathbb{R}^p} \sigma(w^\top x)\; d\nu(w).
\]
The \emph{mass} of the network is the total variation
\[
  |\nu|(\mathbb{R}^p) \;=\; \nu_+(\mathbb{R}^p) + \nu_-(\mathbb{R}^p),
\]
where $\nu = \nu_+ - \nu_-$ is the Jordan decomposition of $\nu$.
\end{definition}

\begin{remark}
  A standard finite-width network $x \mapsto \sum_{j=1}^m a_j \sigma(w_j^\top x)$ is
  recovered as a special case by taking $\nu = \sum_{j=1}^m a_j \delta_{w_j}$, with
  mass $\sum_j |a_j|$.  The infinite-width model strictly generalizes this, allowing
  continuous superpositions over weight space.
\end{remark}

\subsection{Univariate integral representation (Proposition 3.1)}

The following proposition shows that any differentiable univariate function (shifted to
vanish at 0) is \emph{exactly} an infinite-width threshold network, with the signed
measure given by $d\nu = g'(b)\, db$.

\begin{proposition}[Univariate integral representation]\label{prop:univariateIntegralRep}
  \lean{Approximation.InfiniteWidth.univariateIntegralRep}
  \uses{def:infiniteWidthNetwork, def:thresholdActivation}
Let $g : \mathbb{R} \to \mathbb{R}$ be differentiable on $[0,1]$ with $g(0) = 0$ and
$g' \in L^1([0,1])$.  Then for all $x \in [0,1]$,
\[
  g(x) \;=\; \int_0^1 \mathbf{1}[x \ge b]\, g'(b)\; db.
\]
\end{proposition}

\begin{proof}
  \uses{prop:univariateIntegralRep}
By the fundamental theorem of calculus and $g(0) = 0$,
\[
  g(x) = g(0) + \int_0^x g'(b)\, db = \int_0^1 \mathbf{1}[x \ge b]\, g'(b)\, db. \qedhere
\]
\end{proof}

\begin{remark}
  The cost of approximating $g$ by sampling from this infinite-width representation
  scales with $\int_0^1 |g'(b)|\, db$ (the total variation of $g$), which is
  \emph{adaptive}: flat regions of $g$ contribute nothing.
  This contrasts with the grid-based approximation of Proposition~2.1, which pays for
  every grid cell regardless of local flatness.
\end{remark}

% -------------------------------------------------------------------------
\section{Barron norm and Fourier representation}
\label{sec:barron-norm}

\subsection{The Barron norm}

\begin{definition}[Barron norm]\label{def:barronNorm}
  \lean{Approximation.BarronNorm.barronNorm}
The \emph{Barron norm} of a function $f : \mathbb{R}^d \to \mathbb{R}$ is
\[
  \|f\|_{\mathrm{Barron}}
  \;:=\; \int_{\mathbb{R}^d} \bigl\|\widehat{\nabla f}(w)\bigr\|\; dw
  \;=\; 2\pi \int_{\mathbb{R}^d} \|w\| \cdot |\widehat{f}(w)|\; dw,
\]
where the identity follows from the Fourier derivative formula
$\widehat{\nabla f}(w) = 2\pi i\, w\, \widehat{f}(w)$.
\end{definition}

\begin{definition}[Barron class]\label{def:barronClass}
  \lean{Approximation.BarronNorm.BarronClass}
  \uses{def:barronNorm}
The \emph{Barron class with norm $C$} is
\[
  \mathcal{F}_C \;:=\; \bigl\{f : \mathbb{R}^d \to \mathbb{R}
    \;\big|\; \widehat{f} \in L^1,\; \|f\|_{\mathrm{Barron}} \le C\bigr\}.
\]
\end{definition}

\begin{remark}\label{rem:barronNormDimensionFree}
  For functions in $\mathcal{F}_C$, Barron's theorem (below) shows that the number of
  nodes required to achieve $L^2(P)$ approximation error $\varepsilon$ is $O(C^2/\varepsilon^2)$,
  independent of the input dimension $d$.
  This is in stark contrast to the folklore constructions of Chapter~2, which require
  exponentially many nodes.
\end{remark}

\subsection{Examples of finite Barron norm}

\begin{proposition}[Gaussian Barron norm]\label{prop:gaussianBarronNorm}
  \lean{Approximation.BarronNorm.gaussian_barronNorm_bound}
  \uses{def:barronNorm}
Let $f(x) = (2\pi\sigma^2)^{d/2}\exp(-\|x\|^2/(2\sigma^2))$ be the unnormalized Gaussian
with variance $\sigma^2$.  If $2\pi\sigma^2 \ge 1$, then
\[
  \|f\|_{\mathrm{Barron}}
  \;\le\; \frac{\sqrt{d}}{\sqrt{2\pi}\,(2\pi\sigma^2)^{(d+1)/2}}.
\]
In particular, $\|f\|_{\mathrm{Barron}} = O(\sqrt{d})$ when $\sigma^2 = \Theta(1/(2\pi))$.
\end{proposition}

\begin{proof}
  \uses{def:barronNorm, prop:gaussianBarronNorm}
The Fourier transform of the Gaussian is
$\widehat{f}(w) = \exp(-2\pi^2 \sigma^2 \|w\|^2)$,
which is an unnormalized Gaussian with variance $(4\pi^2 \sigma^2)^{-1}$.
Applying Hölder's inequality to $\int \|w\| |\widehat{f}(w)| dw$ against the normalizing
constant of $\widehat{f}$ gives the stated bound; see Barron~(1993), Section~IX.9.
\end{proof}

\subsection{Barron's representation theorem}

The main result of Barron~(1993) expresses any function with finite Barron norm as an
exact infinite-width representation.

\begin{theorem}[Barron representation]\label{thm:barronRepresentation}
  \lean{Approximation.BarronNorm.barronTheorem}
  \uses{def:barronNorm, def:infiniteWidthNetwork, def:thresholdActivation}
Suppose $f : \mathbb{R}^d \to \mathbb{R}$ satisfies $f \in L^1$, $\widehat{f} \in L^1$,
and $\|f\|_{\mathrm{Barron}} < \infty$.
Write $\widehat{f}(w) = |\widehat{f}(w)| \exp(2\pi i\,\theta(w))$.
Then for all $\|x\| \le 1$,
\begin{align*}
  f(x) - f(0)
  &= \int_{\mathbb{R}^d}
    \frac{\cos(2\pi w^\top x + 2\pi\theta(w)) - \cos(2\pi\theta(w))}{2\pi\|w\|}
    \cdot \bigl\|\widehat{\nabla f}(w)\bigr\|\; dw \\
  &= -2\pi\iint_0^{\|w\|}
    \mathbf{1}[w^\top x - b \ge 0]
    \bigl[\sin(2\pi b + 2\pi\theta(w)) |\widehat{f}(w)|\bigr]\; db\, dw \\
  &\quad + 2\pi\iint_{-\|w\|}^0
    \mathbf{1}[-w^\top x + b \ge 0]
    \bigl[\sin(2\pi b + 2\pi\theta(w)) |\widehat{f}(w)|\bigr]\; db\, dw.
\end{align*}
The last two lines express $f(x) - f(0)$ as an infinite-width threshold network whose
corresponding signed measure $\mu$ satisfies $|\mu|(\mathbb{R}^{d+1}) \le 2\|f\|_{\mathrm{Barron}}$.
\end{theorem}

\begin{proof}
  \uses{thm:barronRepresentation, def:barronNorm}
The first equality follows from the Fourier inversion formula and the polar decomposition
of $\widehat{f}$, using the identity
\[
  \cos(2\pi w^\top x + 2\pi\theta) - \cos(2\pi\theta)
  = \int_0^{w^\top x} -2\pi \sin(2\pi b + 2\pi\theta)\, db.
\]
The threshold representation in the second and third lines converts the integral over $b$
from $0$ to $w^\top x$ into indicator functions $\mathbf{1}[\cdot]$.
The mass bound $\le 2\|f\|_{\mathrm{Barron}}$ follows by estimating
$|\sin(\cdot)| \le 1$ and integrating over $b \in [-\|w\|, \|w\|]$.
\end{proof}

% -------------------------------------------------------------------------
\section{Sampling from infinite-width networks}
\label{sec:sampling}

\subsection{Maurey's lemma in Hilbert spaces}

The following lemma underpins all sampling-based approximation results in this chapter.

\begin{lemma}[Maurey sampling; Pisier 1980]\label{lem:maureySampling}
  \lean{Approximation.Sampling.maureySampling}
Let $H$ be a Hilbert space and let $V$ be a random variable supported on $S \subseteq H$
with $X = \mathbb{E}[V]$.  Let $V_1, \ldots, V_k$ be iid draws from the distribution of $V$.
Then
\[
  \mathbb{E}\Bigl[\Bigl\|X - \tfrac{1}{k}\sum_{i=1}^k V_i\Bigr\|^2\Bigr]
  \;\le\; \frac{\mathbb{E}[\|V\|^2]}{k}
  \;\le\; \frac{\sup_{U \in S} \|U\|^2}{k}.
\]
\end{lemma}

\begin{proof}
  \uses{lem:maureySampling}
Write $\frac{1}{k}\sum_i V_i - X = \frac{1}{k}\sum_i (V_i - X)$.  Then:
\begin{align*}
  \mathbb{E}\Bigl[\Bigl\|\tfrac{1}{k}\sum_i (V_i - X)\Bigr\|^2\Bigr]
  &= \frac{1}{k^2} \mathbb{E}\Bigl[\sum_i \|V_i - X\|^2
    + \sum_{i \ne j} \langle V_i - X,\, V_j - X\rangle\Bigr] \\
  &= \frac{1}{k} \mathbb{E}[\|V - X\|^2]
  \;\le\; \frac{\mathbb{E}[\|V\|^2]}{k},
\end{align*}
where the cross terms vanish by independence and zero mean, and the last step uses
$\mathbb{E}[\|V - X\|^2] \le \mathbb{E}[\|V\|^2]$.
\end{proof}

\begin{lemma}[Existence of deterministic sample]\label{lem:maureySamplingExistence}
  \lean{Approximation.Sampling.maureySamplingExistence}
  \uses{lem:maureySampling}
Under the same hypotheses, there exist deterministic $u_1, \ldots, u_k \in S$ with
\[
  \Bigl\|X - \tfrac{1}{k}\sum_{i=1}^k u_i\Bigr\|^2
  \;\le\; \frac{\sup_{U \in S} \|U\|^2}{k}.
\]
\end{lemma}

\begin{proof}
  \uses{lem:maureySamplingExistence, lem:maureySampling}
The probabilistic bound in Lemma~\ref{lem:maureySampling} guarantees that the expectation
is at most the right-hand side; therefore at least one realization $(u_1, \ldots, u_k)$
achieves the bound. (Probabilistic method.)
\end{proof}

\subsection{Maurey's lemma for signed measures}

\begin{lemma}[Maurey for signed measures]\label{lem:maureySignedMeasure}
  \lean{Approximation.Sampling.maureySamplingSignedMeasure}
  \uses{lem:maureySampling, def:infiniteWidthNetwork}
Let $\mu$ be a nonzero signed measure on $S \subseteq \mathbb{R}^p$ with Jordan
decomposition $\mu = \mu_+ - \mu_-$ and total mass
$\|\mu\|_1 = \mu_+(\mathbb{R}^p) + \mu_-(\mathbb{R}^p)$.
Let $P$ be a probability measure on inputs $x$, and write
\[
  g(x) \;:=\; \int_{\mathbb{R}^p} g(x; w)\; d\mu(w).
\]
Define the normalized distribution $\widetilde{\mu}$ by sampling $s \in \{+1, -1\}$ with
$\Pr[s = +1] = \|\mu_+\|_1 / \|\mu\|_1$, then $w \sim \mu_s / \|\mu_s\|_1$,
and output $\widetilde{g}(\cdot; (w, s)) = s\,\|\mu\|_1\, g(\cdot; w)$.

Then for iid draws $(\widetilde{w}_1, \ldots, \widetilde{w}_k) \sim \widetilde{\mu}$,
\[
  \mathbb{E}\Bigl[\Bigl\|g - \tfrac{1}{k}\sum_{i=1}^k \widetilde{g}(\cdot;\widetilde{w}_i)
    \Bigr\|_{L_2(P)}^2\Bigr]
  \;\le\; \frac{\|\mu\|_1^2 \sup_{w \in S} \|g(\cdot; w)\|_{L_2(P)}^2}{k},
\]
and there exist deterministic $(w_1, s_1), \ldots, (w_k, s_k)$ with $w_i \in S$,
$s_i \in \{+1,-1\}$, achieving the same bound.
\end{lemma}

\begin{proof}
  \uses{lem:maureySignedMeasure, lem:maureySampling}
One verifies that $\mathbb{E}_{\widetilde{\mu}}[\widetilde{g}(\cdot;(w,s))] = g(\cdot)$
by unfolding the definition of $\widetilde{\mu}$ and using $\mu = \mu_+ - \mu_-$.
Then apply Lemma~\ref{lem:maureySampling} in the Hilbert space $L_2(P)$ to the random
variable $V = \widetilde{g}(\cdot;\widetilde{w})$ with
$\|V\|_{L_2(P)} \le \|\mu\|_1 \sup_w \|g(\cdot;w)\|_{L_2(P)}$.
The existence of the deterministic sample follows by the probabilistic method
(Lemma~\ref{lem:maureySamplingExistence}).
\end{proof}

\subsection{Barron's sampling bound}

Combining the Barron representation (Theorem~\ref{thm:barronRepresentation}) with
Maurey's lemma for signed measures (Lemma~\ref{lem:maureySignedMeasure}) yields the
main quantitative approximation result.

\begin{corollary}[Barron sampling bound]\label{cor:barronSamplingBound}
  \lean{Approximation.Sampling.barronSamplingBound}
  \uses{thm:barronRepresentation, lem:maureySignedMeasure, def:barronClass}
Let $f \in \mathcal{F}_C$ and let $P$ be a probability measure supported on
$\{x : \|x\| \le 1\}$.
Then for any $k \ge 1$ there exist weights $w_1, \ldots, w_k \in \mathbb{R}^d$,
biases $b_1, \ldots, b_k \in \mathbb{R}$, and signs $s_1, \ldots, s_k \in \{+1,-1\}$
such that the $k$-node threshold network
\[
  \widehat{f}(x) \;:=\; f(0) + \frac{2C}{k} \sum_{i=1}^k s_i\, \mathbf{1}[w_i^\top x \ge b_i]
\]
satisfies
\[
  \|f - \widehat{f}\|_{L_2(P)}^2 \;\le\; \frac{4C^2}{k}.
\]
In particular, to achieve $L_2(P)$ error at most $\varepsilon$, it suffices to take
$k \ge 4C^2/\varepsilon^2$ nodes.
\end{corollary}

\begin{proof}
  \uses{cor:barronSamplingBound, thm:barronRepresentation, lem:maureySignedMeasure}
By Theorem~\ref{thm:barronRepresentation}, $f(x) - f(0)$ equals an infinite-width
threshold representation with signed measure $\mu$ satisfying $\|\mu\|_1 \le 2C$.
Each threshold node $\mathbf{1}[w^\top x \ge b]$ satisfies
$\|\mathbf{1}[w^\top \cdot \ge b]\|_{L_2(P)} \le 1$.
Applying Lemma~\ref{lem:maureySignedMeasure} with these parameters gives the bound
\[
  \|f(\cdot) - f(0) - \tfrac{1}{k}\sum_i \widetilde{g}(\cdot; (w_i, s_i))\|_{L_2(P)}^2
  \;\le\; \frac{(2C)^2 \cdot 1}{k} \;=\; \frac{4C^2}{k}. \qedhere
\]
\end{proof}

\begin{remark}[Comparison with folklore constructions]
  \uses{cor:barronSamplingBound}
The folklore constructions of Chapter~2 require exponentially many nodes in the dimension
$d$ to approximate arbitrary Lipschitz functions.
In contrast, Corollary~\ref{cor:barronSamplingBound} achieves dimension-free approximation
(the $4C^2/k$ bound does not depend on $d$) for the subclass $\mathcal{F}_C$.
The price is that functions with exponentially growing Barron norm — such as general
radial functions, or the worst-case lower bounds in Barron~(1993) — pay the same
exponential cost implicitly through $C$.
\end{remark}

% -------------------------------------------------------------------------
\section{Relationship to the NTK and initialization}
\label{sec:barron-ntk}

\begin{remark}[Barron norm and networks near initialization]\label{rem:barronNTK}
As noted in Remark~3.3 of the lecture notes (and elaborated in Ji, Telgarsky, and Xian 2020),
the Barron norm also controls the complexity of shallow networks in the \emph{Neural Tangent
Kernel (NTK)} regime — that is, networks staying near their random initialization.
This connection arises because threshold units are (informally) the derivatives of ReLU
units, so the Barron representation with threshold nodes maps directly onto the linearized
network $f_0$ that arises in the Taylor expansion near initialization.

Formalizing this connection would require extending the NTK formalism to the LML library;
we leave this as future work.
\end{remark}
